% SPDX-License-Identifier: Apache-2.0
% covers: EthSlots
\datasheet{EthSlots}{The Ethernet slot registers on AXI-Lite: LiteEth's
map, with the frame buffers in main memory.}

\dsfacts{\code{lib/parts/src/ethslots.rs}}%
{\code{txhdl\_parts::ethslots::EthSlots}}{\code{//docs:examples}}

\subsection*{Function}

\code{EthSlots} is the register map of Zephyr's
\code{drivers/ethernet/eth\_litex\_liteeth.c} in hardware, so that
bringing up Ethernet is a port of a driver that already works rather
than a design nobody has driven.

Its AXI-Lite side is one port, \code{bus}, a \code{LitePort} with
32-bit addresses and words: \code{bus\_aw}, \code{bus\_ar} and
\code{bus\_w} in, and \code{bus\_b} and \code{bus\_r} out. Seven
further ports face the engines that move the bytes: \code{tx\_busy},
\code{rx\_busy}, \code{rx\_len} and \code{rx\_which} in, and
\code{tx\_base}, \code{tx\_bytes}, \code{tx\_start} and
\code{rx\_base} out. \code{irq} is high while a received frame is
unacknowledged and the receive interrupt is enabled.

It holds no frame. There are two slots each way, used in turn, and
each is a flat buffer of 2048 bytes in main memory rather than inside
the peripheral, so that \code{LineFetch} and \code{LineStore} move
the bytes instead of the processor. The unit's work is to say which
buffer and how many bytes, and to tell the driver when a frame has
landed.

Two slots a direction rather than a descriptor ring, because Zephyr's
driver model never hands a chain to hardware:
\code{net\_pkt\_read(pkt, buf, len)} and \code{net\_pkt\_write(pkt,
buf, len)} take a flat pointer and a byte count, so the hardware only
ever sees one contiguous region per frame.

\subsection*{Parameters}

\code{BASE}, where the four buffers begin. A slot is \code{BASE + n *
2048}, receive slots first and transmit slots 4096 bytes above them.
The board uses \code{0x4100\_0000}, which \code{isa::ETH\_BUF\_BASE}
states.

The two directions are separate regions and have to be. A slot number
alone would make transmit slot zero and receive slot zero the same
address, and a frame arriving would land on one waiting to go out.

\textbf{\code{BASE} must be 1 KB aligned, and nothing checks it.} A
slot is then 1 KB aligned too, since the slots are 2 KB apart, and a
256 beat burst of words is exactly 1 KB, so no burst can cross AXI4's
4 KB boundary. An unaligned \code{BASE} puts that obligation on the
caller, and a burst that straddles a boundary is a protocol violation
rather than a wrong address, so it will not look like a fault in this
unit.

\subsection*{Register map}

The unit decodes the word from address bits 5 to 2, so the words are
at offsets \code{0x0} to \code{0x3c} of its range. The board puts
them at \code{0x3400}, which \code{isa::ETH\_BASE} states.

\dsregs{EthSlots}

Words 10 to 15 are named by no register and read zero.
\code{tx\_ev\_pending} is never set, as the behaviour below says, and
\code{tx\_ev\_enable} is kept so the driver can write it.

\emph{An unnamed word does not read as zero.} Words 4, 5 and 6, which
are write only, and words 10 to 15, which are nothing, all read as
\code{tx\_ev\_enable}. That is issue 454 and it is a fault rather
than a decision: a driver that writes \code{tx\_length} and reads it
back gets neither its value nor zero but an unrelated bit, so neither
a match nor a mismatch means anything. It reads as zero only while
\code{tx\_ev\_enable} is clear, which is why nothing had noticed.

\emph{The transmit side has no interrupt.} LiteEth's driver disables
it and polls \code{tx\_ready}, because Zephyr's send may block, so
\code{tx\_ev\_pending} is never set by anything and
\code{tx\_ev\_enable} gates nothing. Both words exist because the
driver acknowledges them. A driver that waited on
\code{tx\_ev\_pending} would wait for ever.

\dstables{EthSlots}

\subsection*{Behaviour}

\begin{itemize}
\item A read is answered in the cycle the request is taken, and a
write is answered with \code{bus\_b}. The unit never stalls: it needs
nothing the other AXI-Lite peripherals do not.
\item \textbf{An arrival is the store engine falling idle, not the
frame being received.} \code{LineStore} holds \code{running} high
until the write response comes back from memory, so its fall is the
first moment the frame is certainly there. The unit takes that level
rather than a completion pulse, deliberately: a pulse could be joined
to the receiver by somebody reading the map and not the comment, and
the interrupt would then tell a driver to read a frame whose last
beats are still in flight.
\item On an arrival the unit latches \code{rx\_which} into
\code{rx\_slot} and \code{rx\_len} into \code{rx\_length}, and sets
\code{rx\_ev\_pending}.
\item \textbf{An acknowledgement in the same cycle as an arrival does
not lose the frame.} The drives are written with the acknowledgements
first and the arrival after them, because \code{with!} applies its
entries in order and the last drive of a field wins. The other order
reads more naturally and clears the pending bit that an arrival has
just set, leaving the frame in memory, whole and correct, announced
to nobody. It is the cycle where a busy link acknowledges one frame
as the next one lands.
\item \code{tx\_start} is a request rather than a pulse the engine
must catch. Writing one sets a flag; the unit drives \code{tx\_start}
high while the flag is set and the engine is not busy, and clears the
flag when the engine goes busy. So the request lasts until it is
taken.
\item \code{tx\_base} is \code{BASE + 4096 + tx\_slot * 2048} and
\code{rx\_base} is \code{BASE + rx\_which * 2048}. The receiving side
is told where the slot it is filling begins, and which slot that is
comes from the engine, so that a frame lands somewhere the driver is
not reading.
\item \code{irq} is \code{rx\_ev\_pending} and \code{rx\_ev\_enable}
together.
\end{itemize}

\subsection*{What a driver must do}

\textbf{Read the last word of a frame back before writing
\code{tx\_start}.} Stores on this machine are posted and \code{fence}
orders nothing, which is issue 432, so without that read the fetch
engine may read the slot before the frame has landed and send what
was there before.

What makes the obligation dischargeable is measured rather than
assumed: this link answers a load behind an unanswered store to the
same address with the stored value, which
\code{a\_load\_behind\_an\_unanswered\_store\_to\_one\_address} in
\code{lib/parts/src/bus/axi/sim.rs} asserts, and \code{//docs:axi}
states.

\textbf{That is this interconnect's behaviour and not AXI's
guarantee.} AXI does not order a read behind a write, and a driver
written for another fabric on the strength of this paragraph would be
wrong. The reason the read works here is the link, not the memory
being a serialisation point.

\subsection*{Verification}

\begin{itemize}
\item \code{EthSlots} has no \code{\#[test]} of its own.
\item \code{ex\_ethslots} drives the map as
\code{eth\_litex\_liteeth.c} will, in the order that driver does it:
\code{tx\_ready} one while idle and zero while a frame is going, no
interrupt while the store is in flight and one after it, and the
pending bit tested both ways, since a write of zero must not
acknowledge.
\item It also tests the cycle where an acknowledgement meets an
arrival. The cycle to issue at was measured against both orders of
the drives rather than reasoned about, because two earlier attempts
at that test passed against the faulty order.
\item The build simulates the netlist against that run under nvc and
Verilator: \code{//docs:ethslots\_sim\_ethslots\_tb\_test} and
\code{//docs:ethslots\_vsim\_test}.
\item On the board, with the engines behind it, the two Ethernet tests
in \code{cpu/vreteno/tests/board.rs} run programs on the core that
use its registers as the driver does: one waits for an arrival and
reads the frame from the slot and length it names, and the other
sends two frames back to back, waiting on \code{tx\_ready} between
them.
\end{itemize}

\subsection*{Used by}

The board, at \code{cpu/vreteno/src/board.rs}, on the fifth AXI-Lite
slot at \code{0x3400}, with its arrival as the third source of
\code{Plic3}. Its engine-facing ports go to the fetch engine and
\code{FrameOut} on the way out, and to \code{FrameIn} and the store
engine on the way in (issue 151). The example \code{ex\_ethslots}.

\subsection*{Limits}

It moves no frames itself: the engines behind it do. Without them,
as in \code{ex\_ethslots}, a driver binds and reads and writes every
register, and no frame arrives or leaves.

Two slots a direction is queue depth and not throughput. A ring, if
it is ever wanted, is an addition rather than a rewrite: a descriptor
states the same three values \code{tx\_slot}, \code{tx\_length} and
\code{tx\_start} already state.

The unit checks no more of the address than the word, because the
bridge in front of it sends it only its own range.
